\documentclass[11pt,a4paper]{report}
\usepackage[utf8]{inputenc}
\usepackage[margin=1cm,includehead,includefoot]{geometry}
\usepackage{multicol}
\usepackage{fancyhdr}
\usepackage{titlesec}
\usepackage{graphicx}
\usepackage{pdfpages}
\usepackage{ulem}
\usepackage{dirtytalk}
\usepackage{amsmath}
\usepackage{breqn}
\usepackage{url}
\usepackage{amsthm}
\usepackage{epigraph}
\usepackage[super]{nth}
\usepackage{fbb}
\usepackage{changepage}

\usepackage[T1]{fontenc}

\usepackage{csquotes}
\usepackage{amssymb}
\usepackage{breakurl}
\def\UrlBreaks{\do\/\do-}

\setlength{\parindent}{0em}
\setlength{\parskip}{1.2em}

\hyphenpenalty=8000
\tolerance=1000

\newenvironment{Figure}
{\par\medskip\noindent\minipage{\linewidth}}
{\endminipage\par\medskip}


\titlespacing*{\subsubsection}{0cm}{0cm}{0cm}

\newcommand{\triast}{\bigskip\par\noindent\parbox{\linewidth}{\centering\large{*}\\[-4pt]{*}\hskip 0.75em{*}}\bigskip\par}%
\newcommand{\threeast}{\bigskip\par\centerline{*\,*\,*}\medskip\par}%
\newcommand{\oneast}{\bigskip\par{\large\centerline{*\medskip}}\par}%

\usepackage{enumitem}
\setitemize{noitemsep,topsep=0pt,parsep=0pt,partopsep=0pt}
\setenumerate{noitemsep,topsep=0pt,parsep=0pt,partopsep=0pt}

\titleformat{\chapter}[display]
{\normalfont\bfseries}{}{0pt}{\Huge}

\pagestyle{fancy}
\fancyhf{}
\fancyhead[R]{\leftmark}
\fancyhead[L]{The Librarian}
\fancyfoot[C]{\thepage}

\renewcommand{\footrulewidth}{0.5pt}
\setlength{\headsep}{1.5em}

\newcounter{count}

\begin{document}
	
	\includepdf[pages={1}]{cover.pdf}
	\setlength{\parskip}{0em}
	\tableofcontents
	\setlength{\parskip}{1.25em}
		
\chapter{Library}\label{introduction}

\section{Library news}\label{library-news}

\begin{multicols}{2}
	
The Library now sells both white and red poppies. The Chair, Jonny
Heywood, noted that ``the library primarily exists to increase choice
and availability - choice and availability of books, newspapers and
other media. Something's being provided does not, of course, imply
endorsement. To that end, both white and red poppies are now available
in the library.''

The Peace Pledge Union, who manufacture white poppies, say that ``they
represent remembrance for all victims of war, a commitment to peace and
a challenge to attempts to glamorise or celebrate war.''

The British Legion writes that ``The Legion advocates a specific type of
Remembrance connected to the British Armed Forces, those who were
killed, those who fought with them and alongside them.'' It writes of
white poppies that it ``see{[}s{]} no conflict in wearing the red poppy
alongside the white poppy.''

The Library Committee's Krispy Kreme sale was extremely successful -
over a dozen boxes were sold in under fifteen minutes.


\end{multicols}

\section{Editor's note}\label{editors-note}

\begin{multicols}{2}
	
Hume wrote once:

\begin{displayquote}

In every system of morality, which I have hitherto met with, I have
always remarked, that the author proceeds for some time in the ordinary
way of reasoning, and establishes the being of a God, or makes
observations concerning human affairs; when of a sudden I am surprised
to find, that instead of the usual copulations of propositions, is, and
is not, I meet with no proposition that is not connected with an ought,
or an ought not. This change is imperceptible; but is, however, of the
last consequence. For as this ought, or ought not, expresses some new
relation or affirmation, 'tis necessary that it should be observed and
explained; and at the same time that a reason should be given, for what
seems altogether inconceivable, how this new relation can be a deduction
from others, which are entirely different from it. But as authors do not
commonly use this precaution, I shall presume to recommend it to the
readers; and am persuaded, that this small attention would subvert all
the vulgar systems of morality, and let us see, that the distinction of
vice and virtue is not founded merely on the relations of objects, nor
is perceived by reason.

\end{displayquote}

What Hume meant was that an argument must have at least one moral
premise before it can logically make a claim that an action is wrong
based on factual analysis. An example may clarify this. That an action
would cause deaths does not necessarily mean that it is bad - one must
also claim that death is bad.

\textit{The Librarian} has hitherto been split into a dialectic and
scientific-cum-mathematic section. The latter's name is potentially
inadequate; the humanities make no claim to be sciences, but
nevertheless do not always prescribe - often they merely describe. Hence
the scientific-cum-mathematic section will be renamed ``Is'', and
dialectic ``Ought''.

In the coming weeks, we shall, hopefully, smoothen our workflow, so as
to enable issues to be generated faster and on time. This will include
our weekly \textit{Library News} publication, which has not quite met its
publication schedule to a combination of factors.


\end{multicols}

\subsection{General notice}

\begin{multicols}{2}
	
	\textit{The Librarian} is typeset in \LaTeX, with Scribus used to create the front and back pages, in 11pt.
	
	The Editor would also like to remind readers that he accepts letters, agony aunt questions and game submissions, as well as articles. 

\end{multicols}

\chapter{Ought}\label{dialectic}

\section{The Death of Chinese Pluralism}

\textbf{Some of this article has been published before; now that it has been completed, both parts of the article have been put together.}

\subsection{Introductory}

\begin{multicols}{2}
	
About 100 thousand Hong Kong residents marked June 4 with an annual commemoration of the Tiananmen Square massacre. These commemorations provide an important reminder of the true nature of the Chinese state. They are important on their own, by virtue of the sheer barbarity involved in the massacre. Yet they now are perhaps important for us in a self-interested calculus.

China matters. Almost everyone acknowledges now that it matters, and that whatever rhetoric surrounds its rise is not wholly ``hype'', as it might have been of Japan. It matters not only in the UK, where $7\%$ of our electricity supply will soon be reliant on Chinese goodwill and engineering, or in those parts of the Eurasian landmass which are to be connected by its One Belt One Road initiative, but also in the entire world with universal and undeniable force, when the vast majority of the computing equipment essential to the working of the modern world is manufactured in China.

When such a rise occurs, it is important that the nature of the rising power is analysed, both out of an intrinsic need to know of one's fate as interaction with China changes both in intensity and in balance and out of a need to make informed choices as to those dimensions in interaction with China.

First, three initial examples will help to illustrate the scope and aim of this article.

Yang Shuping, a student at the University of Maryland, said in a graduation speech that she felt the "fresh air of free speech" during her time at university Apart from some embellishment in relation to air pollution in her city of Kunming, it's not clear exactly what was taken issue with.

``We need to open up and embrace all the suggestions from outside world," claims the first speaker. "I would be so piss[ed] off if anyone described my country with destruction\footnote{The audio isn't entirely clear.}," she continues.

A few days before, China held a ``Belt and Road initiative" meeting in Beijing. One notable absence was that of Lee Hsien Loong, Prime Minister of Singapore; Singaporean media speculated that this was due to an incident several months prior, wherein a Terrex armed vehicle was impounded by Hong Kong authorities while returning from a training exercise in Taiwan.

Several months earlier, a mainland legal expert proposed an adjustment in the ratio of foreign judges in Hong Kong, after a foreign judge, David Dufton, sentenced police officers who assaulted protestors in the 2014 Occupy Hong Kong protests to jail.

Although all these incidents appear to be unrelated, they are a symptom of a deeper maliaise which touches the heart of every state and establishment, but does so with particularly problematic consequences in states which have experienced recent political ruptures or are lack conventional accountability mechanisms. China, which suffers acutely from both problems, is run by a leadership both in great need and with few sources of legitimacy.

These travails may seem to be the unimportant rumblings of a faraway Eastern nation. They are not. They will soon, if they do not already, affect you and me as much as they affected Yang and Dufton. When China was introverted, it was easy enough for its political oppression to stay within China; now that it has opened up, it is willing and able to exert pressure on outside actors to change their behaviour, often for the worse. It is therefore important to question what motivates the Chinese state and how it comes to decisions, in order to determine what is likely to occur as a result of this externally as well as internally.

This article will explore the ways in which the Chinese state can achieve legitimacy. It concludes that, in the absence of democratic reform, the Chinese state will either collapse or become increasingly stridently nationalist, associating Chinese identity with so-called ``Chinese values'' of dictatorship and international dominance, which presents a number of problems for other states, both in their treatment of Chinese citizens and in their relations with China.

\end{multicols}

\subsection{The foundations of the Chinese nation state}

\begin{multicols}{2}
	
On 12 February 1912, millennia of tradition in the world’s longest continuously surviving civilisation came to a halt, as the child Emperor abdicated. An interregnum marked by instability, military domination by warlords and misery ceased upon communist victory on 1 October. Outwardly, the revolution took the form of a comprehensive movement away from ideals which dated back to the preceding millennium. This was, too, the case inwardly. Between 750,000 and 1,500,000 million were killed in a campaign of mass and popular depravity.

The democratic aspirations of the Chinese people may have been illusory, limited merely to Western-connected intellectuals. Yet 12 February marked also a time in which the Chinese people threw off the shackles of imperial authority. The speed of change is striking. Nominally one of the world’s most powerful people in 1912, Puyi, the child-emperor of China in the last days of empire, spent a decade in Communist prison camps before retiring to write a pro-Mao memoir-cum-apologia, in which he wrote ``I covered up my crimes in order to protect myself”, in reference to his justifications for cooperation with the Japanese in Manchuria.

Inwardly, however, one pressure remained - the need for legitimacy. The removal of millennia of institutional inertia presented the greatest opportunity for democratic renewal that China has ever faced. Although democratic systems occasionally suffer from legitimacy problems due to low turnout, inherent in democratically elected governments is at least some degree of popular support, and so there is a floor of legitimacy. Dictatorships' floor of legitimacy is such that they only are overthrown when a sufficient proportion of the people feel so incensed by regime incompetence or manevolence that they attempt to overthrow the government.

First, we must ask ourselves where legitimacy is derived from, particularly in a Chinese context.

China is not entirely alien to the rest of the world. Its people care about life, liberty (in a sort of collective sense, although perhaps not individually) and other basic needs. In this sense, Chinese regimes need, like other regimes, to ensure that people living under them are adequately fed, housed and clothed. The failure of the republican leadership to do this during the civil war can be said to have lead to their downfall, because it encouraged support for the communists. This need, however, is not absolute. Many Chinese dynasties limped on while their people languished in poverty, and the communists themselves. In that sense, legitimacy derived from living standards exists in China as it exists universally.

There are also ideological concerns and ways of achieving legitimacy. For thousands of years, the idea of the Mandate of Heaven created loyalty to the emperor, at its basic heart. However, it was sufficiently flexible to allow new emperors or dynasties to take over, when ``Heaven'' decided that the emperor or a dynasty had become too venal and corrupt to deserve its mandate. In this way, it provided a useful tool by which to, somewhat external to material concerns, rule. As a result of its great and increasing age, it managed to survive for thousands of years as the basis of Chinese government; we conclude from this that it was quite successful in its aim. 

The Mandate of Heaven is not the only way by which ideological legitimacy can be created. The communists attempted to replace this with a utopian vision of a Maoist society. The degree to which this was successful is quite difficult to ascertain. On the one hand, there was great enthusiasm in the campaign against ``intellectuals'' and other undesirables during the Cultural Revolution and Great Leap Forward. On the other hand, those same people who probably were around persecuting unfortunates quickly reversed their stance and displayed great flexibility in adapting to first ``Socialism with Chinese characteristics'' and then its true face, ie. some form of state capitalism, with all the rent-making opportunities that implied. Were all these hundreds of millions of people really genuinely convinced that Mao was a great visionary who had all the answers, but simultaneously able to adjust? Or, as is more likely, were they moulded by circumstance, perhaps to some degree motivated by genuine fervour, but really largely simulated by each other? One shout to ``kill 1000 intellectuals'' would likely have been followed by another hundred in a sort of collective inaction problem - what if one is the next to be deemed a reactionary for insufficient revolutionary fervour?

How does this apply to the new regime? First, any new state will have problems in finding popular legitimacy. This is inevitable when revolution necessarily causes all the old institutions of government to, at the very least, lose power or fragment. As a result, many basic services which citizens rely on the government to provide and/or would expect of any government, including at the very least security, are likely not to be provided. This problem was compounded by a deliberate policy of repression of ``intellectuals'', ie. people who had a university education or vaguely had knowledge of how to run things which China needed, like power stations or transportation networks. Given this, the foundation of the Chinese state was not built on material concerns - indeed, it had to be actively shifted away from such concerns because initial administration was so poor.

Second, the Chinese state was deprived of the idea of the Mandate of Heaven. Not only were they not monarchical, which meant that such ideas were already largely impossible, given that the Mandate of Heaven was given to an emperor and his/her descendants, collectively comprising a dynasty, they were also diametrically opposed to the intellectual fabric which gave credence to such ideas, in that they sought to destroy traditionalist ideas of God, religion and emperor. We can therefore discount divine legitimacy as a source of inspiration.

Third, what is left is pure ideology. The sole mechanism by which the communists were able to gain legitimacy was ideological struggle, a mechanism which is hard to sustain in this solitary manner. Not only would it have, at some point, come into conflict with the necessity of preventing state collapse - legitimacy is a mechanism to stay in power, but if external factors mean that even with perceived legitimacy one collapses, it does not help, a reliance on ideological struggle fails because of its inherent idealism, which is incompatible with practical problems causing a failure to achieve utopia.

Precipitately, the foundations upon which the Chinese state were built were entirely precarious, and especially ill-suited to China.

\end{multicols}

\subsection{Normalisation}

\begin{multicols}{2}
	
This situation could not go on, and it did not. Many will be familiar with at least an outline of what happened under Deng Xiaoping. Economic liberalisation, and a decision to cease the continual shooting-oneself-in-the-foot policies which marked Maoism, meant that the state could refocus itself on achieving the first type of legitimacy, ie. that based on material concerns.

This worked, for a while. China lifted 800 million people out of poverty and is now the largest economy in purchasing power terms in the world. Yet 10\% growth is barely sustainable. China has been growing at this rate for nearly four decades. It cannot continue for much longer - growth is now down to below 7\%, and may fall further. Independent sources of information which might tell us without government interference what the real economic situation is in China, but until then we are dependent on the government for such information.

The Chinese leadership, however, recognises that it is very difficult to continue to use ideological fervour to sustain the legitimacy it has created based on economic growth. This puts it in a very dangerous situation, because the measure needed to ensure long term growth and hence long term stability are in conflict with measures which would simulate the economy in the short term in order to stave off loss of legitimacy.

It is difficult to determine whether this is true. Government approval ratings in China are often above 90\%, but even if they are conducted impartially, it is possible and indeed quite likely in a nation which went through the trauma of the Cultural Revolution that these polls are affected by culturally ingrained reminders to toe the party line, both because it is correct in an objective sense, and because such toeing is in one's self-interest.

More easily, one can analyse the measures that the Chinese government is taking to ensure that it retains control. This, as a first step, indicates that it is worried about its control of the country.

\end{multicols}

\subsection{Increasing dictatorship}

\begin{multicols}{2}
	
Liberalisation, and democratic stagnation/reversal, in China, largely took place on three tracks: first, internal party mechanisms of control, second, information control, and third, other civil liberties like freedom of movement and freedom from torture.

After the excesses of Mao, the communists in China recognised that they too would have to face up to the ``bad emperor problem'': with a high level of state control, one can enrich oneself and the country with good policymaking, but should a bad ``emperor'' or ``paramount leader'' as the case may be come to power, the whole of China would face severe difficulties. Mao was their ``bad emperor'', and he taught the communists the virtue of true collective leadership. Although the communists have not become paragons of liberal or democratic virtue, it is notable that the infighting which characterised the Party under Mao significantly abated under the mechanism of collective leadership. The solution of collective leadership works by slowing down the acquisition and exercise of power by its diffusion. In doing so, one has fewer ``bad emperors'', even though progress is slowed. It is indubitably superior to the days of Mao or Stalin, and is superior in that it prevents massive shifts in Party power from one faction to another in the hope of achieving total poewr - instead one vies for positions within the Party and the individual powers they offer.

As for information, after Mao, it was significantly freed up. This reached its peak under Zhao Ziyang, but still remained better throughout all of China than under Mao. Chinese media report on a number of social and environmental issues, including smog, contamination from industrial pollutants, and local government incompetence. Although this coverage is often scaled back or retrospectively deleted after censors become nervous, the Chinese people still remember articles after they are deleted. Similarly, somewhat liberal academics were, to some degree, free to debate, within their universities and their dormitories, with less government interference, compared with Mao-era detention and expulsion to the countryside of all academics. In mnay cases, this lead to a degree of introspection which was not present in Mao-era China. This provided great potential for a gradual and optimised transition to democracy and accountability, for it would not have occurred in a vacuum of knowledge.

It cannot be pretended that China was at any point a state which respected certain basic rights of its citizens, like their bodily autonomy or property rights. At no point did the culture of state impunity which pervaded instituions like the ``People's Armed Police'' abate significantly - they were still free to oppress the people. However, what did change was that they were used less. There was less state violence post-Mao. Fewer properties were seized. Torture was used less, and in some cases the judiciary even admitted that it was wrong. This was real progress - certainly not enough for those who still are victimised by it, but still real.

The 18th National Congress of the Communist Party of China marked the start of a turning point in this slow movement towards democratic legitimacy. It initially appeared to be the start of faster movement. Corruption trials of senior party officials finally went ahead. Chinese officials didn't stop mourners from flocking to the house of Zhao Ziyang, a reformist, to mourn his death. Xi was a new face, representing a new era.

These illusions soon wore off. The corruption trials were really a tool to achieve greater power within the Communist party for Xi and his friends. Indeed, Xi himself is said to be suspiciously wealthy, and he and his family hold a number of properties outside China. Chinese officials probably didn't mind letting people mourn Zhao Ziyang because very few people remember him and they felt that their hold on power was greater than ever before, with their control of the main means of transmission of information, ie. the internet. When people posted pro-Zhao messages on internet message boards, they were deleted, but their posters were not rounded up and detained.

Perhaps a new face, Xi is most certainly not doing anything other than shoring up the dictatorial powers of the party.

Xi has been named the ``Core of the Party'', and many commentators suggest that he has more power than any leader since Deng Xiaoping. Some have even suggested that he may wish to continue after his customary pair of five year terms in positions in the Central Military Commission and Party, though he may resign as President as is constitutionally mandated. In his attempts to corral the paty into submission, it is feared that he may have alienated so many officials that the only way to ensure that he is not arrested just as he has arrested some former officials is to stay in power, and that means that we must wait until his death, not his retirement, to see the back of him.

Information controls have been re-imposed. Liberal academics find themselves under increasing pressures to make universities conform as a tool for communism. No longer are the authorities content to leave them as bastions of ineffective and bourgeois liberalism, never seriously mounting a challenge to Communist leadership - they want them to actively participate in their programme of brainwashing the people.

Repression continues unabated. Dissidents are denied access to their lawyers. An environmental documentary which observers initially thought would be permitted under a new censorship policy of tolerating environmental polemics was soon prohibited. Lawyers of dissidents are occasionally detained. Sometimes their lawyers are detained. The state appoints its own lawyers for them, who are often incompetent and fail to defend them at all - the reason why should be obvious.

This is already bad. It indicates that the government of the most populous government in the world does not even pay lip service to ideals of human dignity and accountability, and that anyone who even suggests it ought to runs the risk of finding their lives turned upside down. Actual observance of these norms is out of the question. But it is also normal. Lots of governments have acted like this, and their peoples have still managed to recover to some degree. Many dictatorships have also ruled over large numbers of people. Part I is a reminder that a powerful nation is increasingly at odds with the values which not only the West but much of the world holds dear. Part II of this series will explore the ways in which China is spreading its influence across the world, and why China therefore presents a problem perhaps unprecedented in the history of humanity in its propagation of its own values, if they even exist as such and are not simply base interests.

The broad thesis of the first of this series of articles is that China
is unlikely to liberalise in the near future, and that the most likely
path is one of gradually increasing oppression. This article will
explain why China is increasingly able to implement such goals, both at
home and abroad.

A broad effect of modernity is that everything has become bigger in
magnitude: people, numbers of people, armies, weapons, speeds, data and
coercive capacity have all massively increased in scale, complexity and
efficacy. Hegemony in the tenth century was local, and necessarily so.
Hegemony with modernity affects nearly everyone.

The events predicted in this article are likely to affect vast numbers
of people outside China. They will not necessarily happen instantly, and
there may be reversals within the broad trends described. Nevertheless
the author has a relatively high degree of confidence in the general
direction described.

\end{multicols}

\subsection{Domestic control}\label{domestic-control}

\subsubsection{Propaganda}\label{propaganda}

\begin{multicols}{2}
	
Control of information may not change individual axiology but it
certainly changes individuals' decisions as to how best to implement
that axiology. We should not deny the former, for many of our
instrumental beliefs, and indeed core axiological beliefs, are imparted
by the outside world, and so are influenced by those who control access
to information. Nevertheless, the primary concern of the modern day
Chinese censor is not to change axiological beliefs in the value of
family, riches or success. Rather, it is to change the instrumental
values which are perceived to be the best way to achieve them. Xinhua's
press release on Xi Jinping thought claims that such thought is an
attempt to become "a great modern socialist country that is prosperous,
strong, democratic, culturally advanced, harmonious, and beautiful by
the middle of the century."\footnote{http://news.xinhuanet.com/english/2017-10/24/c\_136702802.htm}

Propaganda's value is derived from its ability to change thought.
Previously, the Chinese censor was constrained by a need to retain some
credibility. The efficacy of Chinese propaganda was necessarily limited
by the tradeoff between message and credibility.

China, however, has moved on from real life. The rise of a peculiarly
Chinese internet has been the primary cause of what has happened. It is
unique in two ways: first, the quality and quantity of its adoption, and
second, the degree of state control over it.

\end{multicols}

\subsubsection{Information control}\label{information-control}

\begin{multicols}{2}
	
China has the world's most internet users - 732 million\footnote{http://www.scmp.com/tech/china-tech/article/2064396/chinas-internet-users-grew-2016-size-ukraines-population-731-million}.
That China's total population is 1.4 billion masks something important -
almost all young people use the internet. The quality of China's
internet usage reflects two important characteristics of this growth.
First, its magnitude incentivises internet development in a way which
smaller markets elsewhere do not. Second, that most of this growth was
recent meant that the way that the internet is used in China is
primarily mobile-led and significantly more adapted to the latest
technologies - epayments and so on - than other markets.

Chinese internet usage is therefore different to that in the West. It
includes much more than its Western equivalents - payment systems for
not only online shopping but also physical shopping, a greater embrace
of the ``gig'' economy in providing platforms for freelancers to sell
their services, and heavy integration between and within apps are
particularly salient features of one representative and popular app -
WeChat\footnote{https://www.economist.com/news/business/21703428-chinas-wechat-shows-way-social-medias-future-wechats-world}.
\emph{The Economist} claims that the unique characteristics of the
Chinese market - its size, a propensity to possess multiple devices and
the cultural phenomenon of red packets - were part of this shift.
Another important factor is, however, that the Chinese internet is
deeply separated from the rest of the world. Although one can reach the
outside world from the inside, and outsiders can reach the Chinese
internet, such connexions are often slow - a number of senior Chinese
political figures have complained about this, as have
scientists\footnote{https://www.hongkongfp.com/2017/03/19/chinese-scientists-speak-great-firewall/}.
The connexion between China and the rest of the world is akin to having
a dirt path between two major cities - though there may be few
insurmountable obstacles, most would not attempt to use it. Users are a
little like drivers, in that they will quickly abandon that which is
inconvenient\footnote{https://www.thinkwithgoogle.com/data-gallery/detail/mobile-site-abandonment-three-second-load/}.
That internet users are fickle means that they are very likely to use
Chinese tools instead - especially since Western tools now do not appear
to offer any particular benefit over their Chinese equivalents.

Chinese control over the internet broadly sues the same technologies
which are available to Western countries - filtering, vast numbers of
administrators and moderators, and other heuristics-based algorithms
which block terms. Two things differ: first, political will, and second,
the overtness with which the Chinese government controls the internet.
In the West, perhaps the most draconian régime is that of the United
Kingdom, where § 58 of the Terrorism Act (2000) states:

\begin{enumerate}
	\item
	A person commits an offence if---
	
	\begin{enumerate}
		\item
		he collects or makes a record of information of a kind likely to be
		useful to a person committing or preparing an act of terrorism, or
		\item
		he possesses a document or record containing information of that
		kind.
	\end{enumerate}
	\item
	In this section ``record'' includes a photographic or electronic
	record.
	\item
	It is a defence for a person charged with an offence under this
	section to prove that he had a reasonable excuse for his action or
	possession.
	\item
	A person guilty of an offence under this section shall be liable---
	
	\begin{enumerate}
		\item
		on conviction on indictment, to imprisonment for a term not
		exceeding 10 years, to a fine or to both, or
		\item
		on summary conviction, to imprisonment for a term not exceeding six
		months, to a fine not exceeding the statutory maximum or to both.
	\end{enumerate}
\end{enumerate}

We should note, of course, that a number of potentially ``reasonable''
excuses have been rejected by the courts, including attempting to
understand terrorist ideology\footnote{https://www.theguardian.com/world/2012/dec/06/woman-jailed-al-qaida-material-on-phone}.

In China, control goes much further. Not only is the scale of censorship
greater, but that there is censorship is censored. Not only are
``terrorist'' materials proscribed - those advocating democracy or
attacking the régime in any way whatsoever are prohibited\footnote{https://www.economist.com/news/special-report/21574631-chinese-screening-online-material-abroad-becoming-ever-more-sophisticated}.

Willingness to politicise the internet - to use its technical
capabilities in a political fashion as opposed to using its technical
capabilities apolitically to facilitate political actions, combined with
its pervasive influence over Chinese lives, is

Any claim that China will liberalise because its people demand it must
overcome this new reality of information control in China. The Communist
Party are now here to stay - they have always had a desire to, and now
they have the power to.

\end{multicols}

\subsection{International control}\label{international-control}

\begin{multicols}{2}
	
China is increasingly able to assert its power abroad. This occurs in
two ways. First, China increasingly has a strategic advantage which
enables it to assert its power over other nations. Second, Chinese
people, in occupying positions elsewhere, are often vulnerable to
mainland influence in a number of ways - cultural, political, familial
and economic.

\end{multicols}

\subsubsection{China's strategic
	advantage}\label{chinas-strategic-advantage}

\begin{multicols}{2}
	
China has, broadly, five advantages, primarily over the United States,
but also over the rest of the world.

First, China has the world's largest manufacturing base\footnote{https://data.worldbank.org/indicator/NV.IND.MANF.CD?locations=CN-US-JP}.
Its advantage here is especially pronounced in the field of electronics.
Informatic systems are important in all fields of life - most Western
economies rely on them to function. No other country is able to match
China in quantity, and, increasingly, quality. China will be able to
exploit this to its advantage as it increasingly becomes less dependent
on other nations for development, as exports as a proportion of its GDP
have decreased\footnote{https://data.worldbank.org/indicator/NE.EXP.GNFS.ZS?locations=CN}
and are likely to decrease further.

Second, China controls the natural resources critical for such
manufacture. China supplies 85\% of the world's rare earths\footnote{https://thediplomat.com/2017/08/revisiting-rare-earths-the-ongoing-efforts-to-challenge-chinas-monopoly/}.
These are crucial in building informatic systems.

Third, China has a population advantage, especially when compared to
other potential competitors. Russia and the United States have
significantly smaller populations at a global level. At a local level,
only India can match its population; all other nations which China will
seek to influence in the coming years - Japan, South Korea, and those in
South East Asia, have far smaller populations.

Fourth, China has a capacity for coördinated action which few other
nations possess. It is relatively trivial to see that American
dysfunction is unlikely to disappear in the near future. India, too,
despite BJP political hegemony, retains its everlasting capacity to slow
reforms down. This gives it an advantage over other nations in adopting
new technology, upgrading its military, increasing economic growth with
infrastructure projects, maintaining domestic political control and
increasing domestic production for strategic purposes.

Fifth, China increasingly leads the world technologically. Consortia of
Chinese scientists have been particularly important in recent
developments in quantum physics\footnote{http://www.bbc.co.uk/news/science-environment-40294795}.
They have also been at the forefront of genetic research\footnote{http://www.nature.com/news/chinese-scientists-to-pioneer-first-human-crispr-trial-1.20302}.

\end{multicols}

\subsubsection{People}\label{people}

\begin{multicols}{2}

There are 50 million Chinese abroad\footnote{http://www.asiapacific.ca/sites/default/files/filefield/researchreportv7.pdf}.
Many of these Chinese occupy important positions in academia, politics,
law and so on. They are particularly prominent in Singapore and South
East Asia as a whole, where their traditional prosperity has often made
them targets for political oppression, as occurred in Indonesia under
Suharto.

Most Chinese have relatives in China. China has no compulsion against
attacking relatives\footnote{http://www.telegraph.co.uk/news/2016/03/28/another-chinese-dissident-says-relatives-detained-over-letter-cr/}\textsuperscript{,}
\footnote{http://www.reuters.com/investigates/special-report/china-uighur/}.
Given this, it appears likely that China will start to, if it has not
already done so, use its overseas Chinese in dubious ways.

This presents an important moral problem. Many Chinese live abroad. It
would normally be abhorrent to impose a blanket ban on the employment of
Chinese in sensitive positions. Nevertheless, when the assumption that
would justify such a ban - that Chinese people cannot be trusted - is
undermined (in this case, by the Chinese government), it is unclear
which should be prioritised. We should expect to increasingly be faced
with such problems.

Chinese people in public life who are not spies will also face problems.
There will be a permanent cloud of suspicion over their heads.

\end{multicols}


\subsubsection{A collective action
	problem}\label{a-collective-action-problem}

\begin{multicols}{2}
	
It is clear that China has a vast number of carrots to offer other
nations in attempting to impose its will. There is, of course, a classic
collective action problem here. Those who refuse to coöperate on a
principled basis lose, but the Chinese people gain nothing from such
refusals, for China will always find someone else to interact with.

Coördinated opposition could come from an alliance of relatively large

\end{multicols}

\subsubsection{Winds of change}\label{winds-of-change}

\begin{multicols}{2}

We underestimate the power that a sense that ``progress'' is one way or
another has. During the end of history, we saw significant
liberalisation and democratisation. This correlated significantly with
two phenomena: first, the fall of the Soviet Union and hangers-on, and
second, the victory of the idea of liberal democracy.

As the Soviet Union fell, dictatorships around the world loosened
controls or fell. In Russia, there was, for a period, genuine freedom of
expression, and, to some degree, liberty. There is evidently none now.
Empirically, there is a correlation.

Why? All dictatorships need a raison d'être, even if merely for show.
This is why dictatorships which have complete control or near complete
control over their societies - the DPRK, China and Russia, for example -
still require some explanation for their actions. No dictatorship has
yet openly declared itself a kleptocracy. This is for two reasons.

First, officials, when overly kleptocratic, cause state collapse, by
overly extracting rent. This can be solved with an ideological basis -
it ensures lip service, allows enforcement mechanisms to be
significantly more effective, and motivates officials in a way which
overt kleptocracy does not.

Second, there is an innate human desire to justify one's actions. It is
possible that humans are incompatible with overt kleptocracy. Even some
of the most evil people that the world has ever seen - Hitler, Stalin
and Mao, attempted to frame their actions as conducive to a certain
(often warped) perception of ``the good''.

It is most certainly not for want of material control that this occurs.

China's dominance means that the winds of progress are perceived to blow
towards it. In the West, confidence in democracy has declined\footnote{https://www.weforum.org/agenda/2016/12/charts-that-show-young-people-losing-faith-in-democracy/}.
Outside it, dictators who previously claimed that the path to democracy
was simply to be had later\footnote{https://www.thetimes.co.uk/article/iraqis-not-ready-for-democracy-says-blairs-envoy-qkf8rm6j3wf}.
Now, it has grown - dictators no longer need to pay lip service to it,
and they now feel free to entirely reject, down to an axiological level,
democracy.

This is important because even though lip service does not always result
in significant gains in freedom for the majority of the population, the
groups which it allows to exist are important in any post-dictatorship
scenario. That, for example, there was continuity between colonial and
postcolonial institutions in many nations appears to have contributed to
stability in a number of British colonies. Where there was no continuity
or this continuity was rejected, instability reigned. We may, for
example, compare the Democratic Republic of the Congo and Namibia. This
is not exactly a fair comparison for a number of reasons, but
empirically nevertheless there is something to be said for continuity.

When dictators feel free to destroy civil society, there is nothing
left. Societies may genuinely be better off with dictators than under
anarchy. Those who would resist do not.

Resistance to oppression involves another collective action problem. If
every single North Korean not part of the régime were to attempt to
overthrow it, the régime would soon fail. Yet this does not happen. For
every individual, the price of resistance is very great, and the likely
gain very limited. Once the collective action problem is overcome,
whether by chance or particularly egregious oppression, there is no more
problem. Civil society organisations are crucial in overcoming
collective action problems. They allow resistance to régimes to
crystallise, providing a focal point for others to join and a core
organisation.

China's ascendency empowers dictators to copy its model, even if it does
not take an active interest.

\end{multicols}

\subsection{Conclusion}\label{conclusion}

\begin{multicols}{2}

China's present leadership has the tools to stay in power for a very
long time, both internally, with its control of information, and
externally, with its strategic advantages. As demonstrated in the first
part of this article, it is not a particularly pleasant one. Its model
of dictatorship will be exported, not primarily deliberately but simply
by its changing of the international political climate. As it does,
other nations will also be able to acquire the tools to entrench their
leaders just as China has done.

\end{multicols}

\chapter{Is}

\section{Adventures in recreational mathematics V: graph theory}

\begin{adjustwidth}{5cm}{5cm}

It has been said that all maths is graph theory.

This is a slight exaggeration. However, graph theory has proven to be an extremely useful branch of mathematics for computer science (in which it is used for modelling social media, travel (e.g. GPS algorithms), computer chip design, and a multitude of other areas), linguistics (in modelling of natural language syntax and word meaning), chemistry, physics, neuroscience, sociology, wildlife conservation, and even urban planning \cite{acityisnotatree}, to name but a few.

\paragraph{Graphs}

Unsurprisingly, graph theory is the study of \textit{graphs}, which bear no relation to bar charts and other statistical diagrams. The most basic kind of graph (a \textit{simple}, \textit{undirected} graph) is defined rigorously as an ordered pair \(G=(V,E)\), where \(V\) is a set of \textit{vertices} and \(E\) a set of \textit{edges}, each of which is a subset of \(V\) containing two elements. This is really just a way of representing a collection of points (which we call vertices (or sometimes points or nodes)) with lines (which we call edges (or sometimes arcs or lines)) drawn between some of them.

\begin{figure}[h]
	\centering
	\includegraphics[width=0.5\textwidth]{simplegraph}
	\caption{A simple, undirected graph}
	\label{simplegraph}
\end{figure}

Figure \ref{simplegraph} would have \(V={A,B,C,D}\), as these are the vertices of the graph, and \(E=\{\{A,B\},\{A,C\},\{B,C\},\{C,D\}\}\), as each of the pairs represents one edge on the graph. We say that a graph is \textit{undirected} if it is not \textit{directed}\footnote{This term will be defined later.\label{definedlater}}.

We say a graph is \textit{simple} if there are no edges which go from one vertex back to the same vertex (such edges are called \textit{loops}) and at most one edge goes between any pair of vertices. Graphs that are not necessarily simple are sometimes called \textit{multigraphs}, \textit{pseudographs}, or \textit{general graphs}, but this term is rather poorly defined, as there is no consensus as to whether these terms allow loops, among other issues; if one uses these terms, one ought to define its meaning before doing so. There is also confusion as to whether \textit{graphs} are by default simple or not; one ought to specify this when this distinction is of any importance.

A graph \(G\)'s \textit{vertex set} (the set of vertices) is notated \(V(G)\), and similarly its \textit{edge set} \(E(G)\). We often abbreviate an edge \(\{A,B\}\) to just \(AB\).

\paragraph{Digraphs}

Digraphs are \textit{directed graphs} (hence the name \textit{digraph}); their primary characteristic that sets them apart from other graphs is that the edges have a direction, thus:

\begin{figure}[h]
	\centering
	\includegraphics[width=0.5\textwidth]{directedgraph}
	\caption{A directed graph}
	\label{directedgraph}
\end{figure}

A digraph is defined rigorously as an ordered pair \(G=(V,E)\), where \(V\) is a set of \textit{vertices} and \(E\) a set of \textit{edges}, just as a normal graph; however, \(E\) instead contains ordered pairs, so that a pair \((X,Y)\) indicates that there is an edge from \(X\) to \(Y\). In Figure \ref{directedgraph}, \(V={A,B,C,D}\) and \(E=\{(A,B),(A,C),(B,C),(C,A),(C,D)\}\). Note that this digraph is still considered simple, as while we have two edges between \(A\) and \(C\), they are in different directions.

Unless otherwise stated, the word \textit{graph} does not usually include digraphs.

\paragraph{Walks}

Much of graph theory concerns itself with \textit{walks} between vertices. A walk from vertex \(X\) to vertex \(Y\) is essentially a way to go from from \(X\) to \(Y\) \say{walking} only along edges; on digraphs, one must \say{walk} in the direction of an edge. For example, in Figure \ref{simplegraph}, \(B \rightarrow A \rightarrow C \rightarrow D\) is a walk from \(B\) to \(D\); however, this is not a valid walk on the directed graph in Figure \ref{directedgraph}, as there is no edge from \(B\) to \(A\). Instead, one must take the walk \(B \rightarrow C \rightarrow D\).

\textit{Paths} are walks where each vertex occurs at most once in the walk. \textit{Trails} are walks which do not go over any edges more than once. \textit{Cycles} are walks that go from a vertex back to itself (they are considered to visit this vertex only once).

A \textit{Hamiltonian path} on a given graph is a path that goes through every vertex exactly once, and \textit{Hamiltonian cycles} are cycles that visit every vertex exactly once. Graphs with Hamiltonian paths are called \textit{Hamiltonian graphs}.

An \textit{Eulerian trail} (or sometimes \textit{Eulerian path}, although this is a misuse of the term path) is a trail that goes through every edge exactly once, and similarly an \textit{Eulerian cycle} is an Eulerian trail that is also a cycle. \textit{Eulerian graphs}, somewhat predictably, are graphs with Eulerian trails.

In fact, the first ever graph theory theorem was about \textit{Eulerian paths}. The famous mathematician Leonhard Euler wanted to know if there was a way to walk through the Prussian city of K\"onigsberg (now Kaliningrad in Russia) and cross each of its seven bridges exactly once (without crossing the river by any other means). K\"onigsberg had two large islands, and the bridges were situated as in Figure \ref{Konigsberg}.

\begin{figure}[h]
	\centering
	\includegraphics[width=0.5\textwidth]{Konigsberg}
	\caption{The Seven Bridges of K\"onigsberg, as they were in Euler's time}
	\label{Konigsberg}
\end{figure}

We can consider this as a graph (shown in Figure \ref{Konigsberggraph}), with the vertices being the land masses and the edges being the bridges; then Euler's walk across all the bridges is an Eulerian trail on our graph.

\begin{figure}[h]
	\centering
	\includegraphics[width=0.5\textwidth]{Konigsberggraph}
	\caption{A graph theoretic representation of the Seven Bridges}
	\label{Konigsberggraph}
\end{figure}

Except for the first and last times we cross an edge/bridge, whenever we enter a vertex/land mass, we then leave it. This means that all vertices/land masses except those at the ends of our trail must be entered as many times as they are left, and so must have even degree.

If we start and end at different vertices/land masses, they must have odd degree, as we leave the former one more time than we enter it, and enter the latter one more time than we leave it. If we start and end at the same vertex/land mass, it has even degree, as we enter it as many times as we leave it.

However, the graph of K\"onigsberg has four vertices, all of which have odd degree. Thus we cannot find an Eulerian trail, as at most two vertices can have odd degree for an Eulerian trail to exist.

It was stated by Euler, and later proven by Carl Hierholzer just before his death in 1871, that we can find an Eulerian trail if and only if there are exactly zero or two vertices of odd degree in a graph, and the graph is \textit{connected}. (A graph is said to be connected if there is a walk between any two vertices. Figure \ref{simplegraph} shows a connected graph, while Figure \ref{disconnectedgraph} does not.)

\begin{figure}[h]
	\centering
	\includegraphics[width=0.5\textwidth]{disconnectedgraph}
	\caption{A disconnected graph}
	\label{disconnectedgraph}
\end{figure}

\paragraph{Trees}

\textit{Trees} are graphs where there is exactly one path between any two vertices, as in Figure \ref{treegraph}. An equivalent definition is that trees are graphs that are both \textit{acyclic} (having no cycles) and connected.

\begin{figure}[h]
	\centering
	\includegraphics[width=0.5\textwidth]{treegraph}
	\caption{A tree.}
	\label{treegraph}
\end{figure}

A \textit{forest} is a \textit{disjoint union} (in graph theory, this essentially means \say{the graph you get by putting several graphs next to each other without adding any new edges}) of several trees.

Every tree or forest is bipartite\footnote{\textit{A City is Not a Tree}}, and every \textit{finite} (having finitely many vertices) tree or forest is planar\footnotemark[\ref{definedlater}].

\paragraph{Bipartite Graphs}

Bipartite graphs are graphs where the vertex set (the set of vertices) of the graph can be divided into two \textit{disjoint} (having no elements in common) sets \(T,U\) so that there are no edges between two vertices in \(T\), and no edges between two vertices in \(U\). It follows that all edges are between vertices in different sets.

An equivalent definition is that graphs are bipartite if and only if they do not contain any \textit{odd cycles} (cycles with an odd number of vertices). One can see why this should be true, as we can divide any single \textit{even cycle} into two sets \(T,U\) to fit our condition by simply alternating which set we put points in as we go round the cycle. It is a little more complex to prove.

It follows that all trees are bipartite, and Figure \ref{bipartitetree} shows the graph from Figure \ref{treegraph} divided into sets \(T,U\), with vertices moved to indicate this.

\begin{figure}[h]
	\centering
	\includegraphics[width=0.5\textwidth]{bipartitetree}
	\caption{The tree shown in Figure \ref{treegraph} divided into two sets so as to prove it to be bipartite.}
	\label{bipartitetree}
\end{figure}

Another equivalent definition is that a graph is bipartite if and only if it is \textit{2-colorable}\footnotemark[\ref{definedlater}].

\paragraph{Complete Graphs}

The \textit{complete graph} \(K_n\) is the graph with \(n\) vertices where there is an edge between every two vertices. Figure \ref{k5graph} shows \(K_5\), the complete graph vertices.

\begin{figure}[h]
	\centering
	\includegraphics[width=0.5\textwidth]{k5graph}
	\caption{\(K_5\)}
	\label{k5graph}
\end{figure}

The \textit{complete bipartite graph} \(K_{m,n}\) is the bipartite graph where set \(T\) has \(m\) vertices, set \(U\) has \(n\) vertices, and there is an edge from every vertex in \(T\) to every vertex in \(U\). Figure \ref{utilitygraph} shows \(K_{3,3}\), which is also known as the utility graph, after the three utilities problem.

\begin{figure}[h]
	\centering
	\includegraphics[width=0.5\linewidth]{utilitygraph}
	\caption{\(K_{3,3}\)}
	\label{utilitygraph}
\end{figure}

\paragraph{Planar Graphs}

The three utilities problem is stated thus:

\begin{displayquote}
	Suppose there are three cottages on a plane (or sphere) and each needs to be connected to the gas, water, and electricity companies. Without using a third dimension or sending any of the connections through another company or cottage, is there a way to make all nine connections without any of the lines crossing each other?
\end{displayquote}

This turns out not to be possible. This is really a graph theory question, and can be phrased more succinctly thus:

\begin{displayquote}
	Is \(K_{3,3}\) a planar graph?
\end{displayquote}

A \textit{planar graph} is a graph that can be drawn on a plane (e.g. a flat piece of paper) so that none of the edges cross each other.

A \textit{subdivision} of a graph \(G\) is a graph which can be obtained by dividing edges of \(G\) and placing new points in the place where we have divided them; we are thus allowed to turn edges into the configuration shown in Figure \ref{subdivision}. More rigorously, we are saying that we are allowed to turn the edge \(\{A,B\}\), for any vertices \(A,B\), into the two edges \(\{A,C\},\{C,B\}\), where \(C\) is a new vertex we have just introduced. We can do this repeatedly, and turn \(\{A,B\}\) into \(\{A,C\},\{C,D\},\{D,B\}\).

\begin{figure}[h]
	\centering
	\includegraphics[width=0.5\textwidth]{subdivision}
	\caption{A subdivision of a single edge.}
	\label{subdivision}
\end{figure}

A \textit{subgraph} of a graph \(G\) is a graph whose vertex set is a subset of that of \(G\), and whose edge set is a subset of that of \(G\). The vertex set of the subgraph must include the endpoints of all the edges of the subgraph, but may also include lone vertices of degree 0.

\textbf{Kuratowski's theorem} states a graph is planar if and only if it does not contain a subgraph that is a subdivision of \(K_5\) or \(K_{3,3}\). This theorem was independently proven by Kazimierz Kuratowski, Orrin Frink, and Paul Smith in 1930. Karl Menger also proved it on graphs where all vertices have degree 3 in 1930.

\paragraph{Colourings}

\textit{Colourings} are ways to label aspects of a graph with colours subject to certain constraints. The two most common types of colourings are vertex colourings, where we colour vertices such that no two vertices with an edge between them are of the same colour, and edge colourings, where we colour edges so that no two edges with an endpoint in common have the same colour. You have probably encountered graph colouring before in Sudoku, which is really a question of colouring an graph with 81 vertices with 9 colours.

The \textit{chromatic number} of a graph is the smallest number of colours required to colour a graph's vertices. We say a graph is \textit{\(n\)-colourable} if its vertices can be coloured with \(\leq{}n\) colours.

The \textbf{four colour theorem} states that any planar graph can be 4-coloured (although it is usually phrased in terms of maps, which, assuming no two separate regions are forced to be the same colour, can be 4-coloured; in fact, map colourings and planar graph colourings are isomorphic).

It was proved in 1976 by Kenneth Appel and Wolfgang Haken, and was the first major theorem proved using a computer. However, the five colour theorem (which says that all planar graphs are 5-colourable) was proved in 1980.

\subsection{Problems}

\begin{enumerate}
	\item Prove the five colour theorem.
	\item A polyhedral graph is a connected graph where the vertices and edges are those of a convex polyhedron (or, equivalently, the removal of any two vertices does not render the graph disconnected). A bicubic graph is both bipartite and cubic. A cubic graph is a graph where all vertices have degree 3. Do all bicubic polyhedral graphs have Hamiltonian cycles?
\end{enumerate}

\subsection{Bibliography}
	[1] 
	Christopher Alexander,
	\textit{A City is Not a Tree},
	1965.

\end{adjustwidth}

\chapter{Review}\label{review}

\section{Lectures}\label{lectures}

This issue's lecture reviews were written by Isky Mathews, Ms.~Stone and
the Editor.

\subsection{Middle East Society - The Great War and its consequences}

\textbf{Professor Ali Ansari, Director of the Institute for Iranian Studies at
the University of St Andrew's}

\begin{multicols}{2}

\subsubsection{Historical context}\label{historic-context}

Before World War I, there were two main imperial bodies in the region
known as the Middle East: Persia and the Ottoman Empire. World War I led
to Ottoman breakup. This was an especially important formative period
for the Middle East. One should always be sceptical of coverage of the
Middle East: Sykes-Picot, ever-present, never saw their plan fully
implemented - nor did anyone else. Russia's importance in the region has
often been underestimated; it has historically had significant interests
in the Middle East.

\subsubsection{The Iraq War}\label{the-iraq-war}

2001 and 2003 disturbed this post-Great War settlement. The Americans
attempted and somewhat managed to reshape the region in their own image.
A particularly salient feature is the innovation of attempting to spread
democracy. Ideally, this would allow a quick turnover; it turned out to
lead to a significant loss of stability.

\subsubsection{Obama}\label{obama}

Obama's legacy is important. All Presidents attempt to make up for their
predecessor's failings. In this case, post-Bush exhaustion led to a
perhaps premature withdrawal from the region. The second main feature of
his rule was a (very) partial attempt to move a little closer to the
pre-1979 situation where Iran was one of the main American partners in
the region. A constructive relationship required a partner; this came
not in the form of Ahmajinedad but Rouhani. The most significant step
forward was nuclear. To some degree, this relationship is not based on
material concerns but was emotional. It was hoped, especially by the
Americans, that initial nuclear progress would lead to a broader
rapprochement. This fits in, perhaps, to belief in Whig history, that
is, history whose moral curve is long but bends towards justice. The
pivotal moment may have been 2011/2012, where sanctions began to cripple
the Iranian economy.

\subsubsection{Modes of interaction}\label{modes-of-interaction}

There are two possible modes of interaction. The first involves the use
of a thin end of the wedge, with the eventual aim of slotting the fat
end of the wedge through too. The second involves having a strong base
of initial progress upon which to build later, in a stable fashion. The
former appears to have been employed under Obama.

Iranian thought values businessmen - Trump fit this description, so the
Iranians thought they could make a deal. They were, however, not
completely correct.

To Trump, the deal's provisions are flawed in one crucial way - they are
limited to a decade. The assumption that after ten years Iran will be
sufficiently sound and integrated into the world economy is at best not
entirely guaranteed. Another limitation of the agreement is that it has
not been legally enshrined - it was pushed through by Presidential
diktat, not legislative consent; just as the then-President gave, so too
the current President may take away.

\subsubsection{Who to trust?}\label{who-to-trust}

Professor Ansari concluded with an anecdote. An analyst, once, noticed
that posters of President Obama had been plastered all over Iranian
cities. He failed to notice that Obama was pictured next to Shemr - an
ancient vilain in Shi'ism.

\subsubsection{Questions}\label{questions}

\begin{enumerate}
	\def\labelenumi{\arabic{enumi}.}
	\item
	How best can the Joint Comprehensive Plan of Action (JCPOA) be
	supported? The JCPOA makes political development contingent on
	economic development. A lack of Western investment in crucial
	industries has blunted its benefits to Iran. Crucially, access to
	dollar markets is lacking in Iran, curtailing investment and growth.
	These two concessions would lead to economic development and goodwill.
	\item
	What is the difference between the two methods of interaction
	previously described? Why must a first step necessarily be unstable?
	The Iranians decided to negotiate exclusively on nuclear issues. All
	further discussion was precluded. Had a grand bargain been struck,
	other issues may have been included; now that nuclear-related and
	sanctions-related issues have been somewhat solved, the Americans have
	fewer bargaining chips for later negotiations or negotiations on the
	opening of negotiations.
\end{enumerate}

\end{multicols}

\subsection{Stipatores Honorati}\label{stipatores-honorati}

\begin{multicols}{2}

After a quick change of venue from the Refectory to Purcell's common
room (the tables and chairs had been cleared out of the Refectory ready
for redecoration over Exeat), cake was served, tea was made and we all
settled down to hear William speak on the history of apple cake. Norfolk
apple buns had been set aside in favour of Dorset apple cake (the buns
do not keep well and are also best made with almonds, which are
obviously a banned substance) but nobody voiced any complaints and the
quick disappearance of everyone's cake spoke well of the quality - your
correspondent, for one, thought it excellent.

William then led us on a quick whirl through various apple cakes from
Dorset and Somerset, a West Country poem and the revelation that Norfolk
apple buns are actually a surprisingly recent invention created by Pye
Baker of Norfolk, finishing off with a reading about apples. Questions
were then welcomed, including why do we so often see apple paired with
cinnamon (no conclusion was drawn - perhaps a topic for further
research) and what everyone's favourite variety of apple was. Many
different suggestions were made, with high scorers being Braeburns and
Cox's, but Asian apples had a late surge and one dissenting voice
plumped for oranges instead.

A spirited debate then sprung up about the differences between muffins
and cupcakes. Some members argued that they were the same thing, but
others highlighted several differences including, but not limited to,
the fact that the muffin escapes its case and the cupcake stays
contained. The meeting then drew to a close with Colin the Caterpillar
being accepted as the November cake, despite not being homemade.

\end{multicols}

\subsection{Biology bites}\label{biology-bites}

\begin{multicols}{2}

This talk at Huxley Society was a rounding up of the science
department's ``Biology Week'', consisting of multiple successive
mini-presentations on different aspects of biological study.

The talk began with Hein Mante's presentation on sharks. He firstly gave
great detail on their evolutionary history:

\begin{itemize}
	\item
	400 million years ago, the first cartilaginous fish have been observed
	to have evolved.
	\item
	At 350 million years ago, shark scales were first formed in various
	pre-shark species.
	\textbf{Placoid scales are, in fact, one of the distinctive features of sharks. They produce small-scale vortices around the sharks’ bodies which improve hydrodynamic efficiency and appear to have evolved from vertebrate teeth, as they contain a central pulp cavity supplied with blood vessels.}
	\item
	320 million years ago, the first shark, \textit{Cladselache}, was
	recognised.
	\textbf{It should be noted that this makes sharks one of the oldest and most successful branch of fish.}
	\item
	Around 100 million years ago, a great explosion in shark species
	occurred for somewhat unknown reasons.
	\item
	65-60 million years ago, the Great White and Megalodon sharks become
	distinct species.
	\item
	1 million years ago, the last Megalodon specimen dies.
\end{itemize}

He then continued on to discuss their predatory behaviour and whether
their ``man-eating'' reputation in public opinion is well-deserved. Hein
mentioned that just 5 deaths per annum are caused by shark attacks, in
comparison to Malaria's 780,000/annum and AIDS' 1,800,000/annum
\textbf{(or the 20 deaths per year caused by improper use of vending machines)}.
He observed that shark size affects predatory behaviour and that 2.3m
appears to be a ``critical length'' beyond which sharks become
significantly more aggressive, attempt to eat much larger prey than
would normally be possible to swallow and also exhibit a much greater
chance of trying to eat indigestible objects, such as pieces of coral,
sponge or human waste.

Hein ended by considering sharks' main predator, orcas. It was mentioned
that they flip sharks upside-down using their tails - it appears this
action causes a shark's brain to flood with endorphins and as a result
they become docile and motionless, even with the orca biting into them!

Mr.~Moore's presentation followed, beginning with an anecdote in which
he had overheard some upper school westminsters declaiming Biology and
those who took it along some remarks suggesting its inferiority to the
other sciences. He jokingly expressed his unhappiness at this comment
and said that Biology is often portrayed as a ``Mickey Mouse'' subject;
he humorously concluded from this that the way to understand this remark
is to examine the history of Mickey Mouse himself.

He noted how Mickey Mouse's appearance and personality has changed
significantly. No longer does Mickey Mouse smoke or steal or commit any
sort of wrong-doing and he has gotten a larger cranium, eyes and
exaggerated limbs. Mr.~Moore explained that humans naturally find these
features cute and pleasurable to look at - in fact, these features are
found in nearly all baby animals and this is part of the reason most
creatures have an urge to look after their young. Konrad Lorenz, the
Nobel Prize winning biologist, was mentioned for his work in describing
and categorising these precise physiological differences and follow up
work by neuroscientist Stephen Hamann showed that the common responses
to these features were involuntary amygdala responses, so nearly all are
lying if they claim to not find such features attractive.

He described the phenomenon known as \textit{neoteny}, where adults of a
species retain features normally only found in juvenile specimens. He
then pointed out that humans appear to have developed a form of neoteny,
since adults of \textit{Homo Sapiens} share many facial characteristics
with young primates.
\textbf{It is believed that this was evolutionarily advantageous because having these attractive features probably made it more likely for a given "neotenic" individual to be helped by others.}

Mr.~Moore concluded that there is Biology in Mickey Mouse but the
opposite is not true.

Brandon Tang then began his short presentation on cannibalism in nature
by showing that it appears \textbf{(perhaps oddly)} often throughout the
animal kingdom: pray mantis who eat their mates, tadpoles eating their
siblings, hippos etc. He detailed a notable example in which a primate
mother kept the body of her dead baby with her for at least a month,
before proceeding to eat increasingly large portions of the body. He
elaborated that this was not uncommon with primates; pack leaders have
been seen to be eaten and families of primates have been seen to travel
as a group eating other primates they encounter.

The traditional explanation for this was that in certain crowded areas
of a forest, there may be great sexual competition and competition for
food and that this leads in certain situations to primates killing
others of their own kind, either out of some form of stress or
necessitation.

Brandon said that he felt these explanations were inadequate as it
doesn't quite explain, for example, why direct relatives display a
tendency to not attack each other but rather group together nor does it
give any explanation for the situation of the mother carrying the dead
baby as mentioned above. He gave a counter proposition: that this
instead suggests such primates have a much higher level of complexity to
their social behaviour than previously thought. He finished with the
conjecture that perhaps their emotional intelligence has been
underestimated and that they are quite capable of complex feelings
towards others.

The final presentation was a talk on sloths and their peculiar habits by
James Ferner. It was mentioned that in Costa Rica, there are two species
of sloths: the three-toed and two-toed sloths. Along with the obvious
difference in the number of claws they have on their hands, the two
species have numerous physiological differences - this observation along
with DNA samples have led to the hypothesis that the two species come
from quite distinct branches of the tree of life.
\textbf{The two-toed sloth is believed to have been tree climbing for some million years but the three-toed sloth may have originated from the Giant Ground Sloth and become tree dwelling much more recently.}
If this was found to be true, James notes, then these sloths species
would demonstrate one of the most striking examples of convergent
evolution ever discovered.

Both species are folivores (leaf-eaters) with 4 stomachs and are
extremely slow-moving. They come down once a month to go to the loo and
this is when they are most vulnerable - they are predated on by snakes,
harpy eagles and cougars. James then mentioned that their hair is so
thick and curly that it is often seen green from the algae cultures that
grow in it. They have developed, it seems, a form of symbiotic
relationship with these plants, since the curls of the hair appear to
exist to trap moisture which promotes algae growth, where the algae's
growth provides a form of camouflage in the foliage to sloths. This
appears to make it more difficult for harpy eagles to find them in the
canopy.

James ended the series of talks by mentioning their mating rituals,
since he believed them to be quite humorous. It appears that females
emit a high-pitch shriek from their tree to make it clear that they are
ready to mate and then male sloths become much more active than normal,
with some two-toed sloths seen swimming across lakes to get to the
female. If multiple males arrive at roughly the same time, they can
fight \textbf{(quite slowly!)} for dominance in mating.

\end{multicols}

\section{American Classics}

\begin{multicols}{2}

The Sixth Form festival is known for producing plays of a very high
standard, taking advantage of talented students both new and old.
\textit{American Classics} continued this trend. Brilliantly entertaining
as well as thought-provoking, all three productions were equally
enjoyable and every cast member performed well. Lightning smooth
transitions and a well-chosen mix of tones make it difficult to find
faults within the play. One criticism of the play as a whole, echoed by
a number of pupils, is its length - at almost three hours long, it was
difficult to comprehend subtler aspects. The three hours, however,
proceeded unflaggingly, without inducing boredom.

The play opened in darkness. Max Raphael narrated each opening scene in
the style of an American circus presenter, with slicked back hair and
the associated voice inflections. We were immediately brought into the
world of American theater. The first scene was one of the opening scenes
from Thorton Wilder's \textit{Our Town} which describes the day to day
lives of the residents of Grover's Corner, a small town in
Massachusetts. However the omniscient narration of this scene is what
provides us with the wider context of the scene, interrupting the
seemingly idyllic town with information about people's deaths, marks the
scene. It was nevertheless the narration in the second part which stood
out. Aryan Mishra's monologue on the eternity of human beings was
beautifully delivered by and acts as in contrast to the first part to
\textit{Our Town}, which marked an ordinary day. The second part showed
one of the characters, recently having died, reminiscing on her past
life, conversing with the dead about the ordinary days that used to
occur. This last scene is extremely powerful, especially considering the
contrast between Mrs. Gibbs in the two scenes, played by Agatha Pethers.
In the first she is cheerful and upbeat, considering selling old
furniture to afford a trip to Paris, which she eventually does for
\$350. In the second part however, she was one of the dead, monotonously
advising Emily Webb on what she should do and not reacting when Dr.
Gibbs places flowers on her grave. This contrast was greatly moving -
coupled with the monologue on human nature, your correspondent found
himself fighting back a single tear, to no avail. As the scene
concluded, the sun rose, ending on a strangely optimistic note full of
ambiguity.

The transitions between each scene were rather striking - smooth and
entertaining, each one very different from the last, ranging from Matt
Carver's shaking of his legs to David Bowie's "Let's Dance" to Theo
Tompkins' hopping surreally through the rest of the cast in slow motion
as a rabbit, they separated each scene and encouraged the audience to
transition from the previous scene.

Unfortunately, one of the few other similarly striking aspects of the
play were the attempts at American accents, and their varying degrees of
success; on the most part they were quite convincing, but there were
times when one was unsure which accent was being attempted, hearing a
strange mix of English and something else that somewhat unidentifiable.
For many, this was not a major issue and did not detract significantly
from their enjoyment of this production.

The play then moved swiftly on, to \textit{Angels in America} by Tony
Kushner, following the life of Prior as he wrestles with HIV in 1980's
New York; in the scene that was chosen, he tells his boyfriend Louis of
his condition. David Edwardes-Ker did an excellent job of conveying the
occasionally hilarious bitterness of Prior, and the almost demented
enthusiasm with which he explains how he is going to die to Louis is
done very well; Louis's (played by Boris Karapetyan) reaction well suits
the audience's similar feelings of discomfort to the news. The second
part, in which a dejected Prior in drag meets a Mormon strung out on
Valium, was probably one of my favorite scenes from the play. Intensely
thought-provoking, both Prior and Harper (played extremely well by
Annika Heath) learn something new from each other in the "threshold of
revelation". Annika gave a hilarious portrayal of Harper and the back
and forth between them was executed very well. The scene ended with a
great monologue by Prior on self identity and worth, and the
realisation, portrayed by David, is suitably emotional and powerful.

Next, in a wonderfully creepy transition involving lots of screaming,
was Arthur Miller's \textit{The Crucible}. The first scene of the two
really captures the panic surrounding the fate of the young girls,
notably Cyci Singh, who showed well the fear that surrounded their
situation, fearing greatly accusations of witchcraft. In a powerful
moment she lashes out against one of her co-conspirators in fear of
being discovered with great anger. Alex Foster as Parris also shows the
fear and nervousness surrounding the affair, fearing that he may be
kicked from his pulpit, with witchcraft having been discovered in his
house. The second part though was very powerful, when Elizabeth Proctor
is arrested for suspicion. Max Raphael as John Proctor stole the show
completely with the immense level of raw emotion he was able to portray
in trying to defend his wife from arrest, and the calm, obliging
response of Rosa Calcraft as Elizabeth was deeply moving. One could hear
a pin drop as he cried out "{[}s{]}he will not go". Yet, the changing of
certain actors for the characters created a little confusion, and it was
hard at times to contextualise events; admittedly, this is a trivial
piece of criticism compared to the rest of the scene.

Next, was the bizarre \textit{Machinal} by Sophie Treadwell, depicting the
life of a nervous, nameless young woman trying to live in the fast pace,
anxiety-inducing world in which she lives. The pace of the first scene
is very well set and quite comedic at times, the repeated phrase "hot
dog" never failed to entertain and one was amazed at the sheer speed at
which the dialogue flew very tightly between the office workers. However
it was the brilliant contrast to this fast pace world, portrayed by
Polly Ruppel, which made the increasingly industrialising world seem
wrong and unnatural; at times the performance was quite chilling.
Especially in the second scene, where she is financially obligated to
marry a man whom she does not love nor feel comfortable around, the
atmosphere is deeply uncomfortable as the scene documents the events of
the honeymoon. The scene became more and more chilling as Jones, whose
self-obsession was wonderfully explored by Leo Doody, puts pressure on
her to do things that she does not want to do; the final piece is her
spiraling into an anxiety attack, heightening the discomfort which had
acted excellently done by Polly.

Tennessee William's extremely entertaining \textit{Glass Menagerie}
followed; it showed Laura hosting a gentleman Caller, Jim, her old high
school crush. Jim's clumsiness and simplicity was perfectly executed by
Theo Tompkins, who by now has mastered the art of the Hugh Laurie-like
glazing of the eyes. The scene unfolds with Laura showing Jim her glass
collection, notably a unicorn which is very precious to her.
Unfortunately during a dance between the two of them, the figure is
knocked off the table, losing its horn. In consoling her Jim kisses her,
regretting his decision, as he is now in a steady relationship. The
painful disappointment of Laura is easily noticeable and well portrayed
by Pandora Mackenzie, in a commentary on identity and individuality.

Next is the pivotal scene from \textit{Who's afraid of Virginia Woolf} in
which George and Martha invite a younger couple back home after a party
and expose their marital problems in a violent way. Beth King was a
hilarious Honey and really made the scene entertaining. Darcy Dixon's
portrayed Martha's taunting manner well, capturing the escalation to
"Violence! Violence!" subtly. Parth Agarwal was also entertaining,
openly hitting on Honey to many laughs from the audience. The anger that
he buildt up in the situation, followed by his sarcastic remark at the
end, was pitch perfect and captured the general hostility well.

Now, the final piece of Act One, Tennessee Williams' \textit{A Streetcar
	Named Desire}. In this scene Blanche Dubois, trying to maintain self
respect, as her life falls around her, has to confront her sister's
brutish husband regarding the loss of some property. Rosa Calcraft makes
an excellent Blanche, with one of the better accents in the play, trying
to maintain self dignity, while trying to evade Stanley's anger through
flirting. Stanley, played suitably directly, by Cyrus Gilmartin, is hard
not to root against, as he accuses and insults Blanche. The scene ends
with the happy news of Stella's pregnancy, and the dynamic between the
two sisters is believable and natural. Stella's concern for Blanche was
shown well to the audience by Annika Heath, and we one really empathised
with her throughout.

Finally, starting the Act Two was \textit{The Children's Hour} wherein a
child who is given what she thinks is an unfair punishment takes her
revenge by naming her two headmistresses as lovers, an accusation which
goes on to destroy their careers. Although being very strong, it was one
of the weaker scenes. A few of the accents were hard to pin point and
the general order and dynamic of dialogue were at times a bit shaky.
That being said, this scene was able to create a great tension between
the grandmother, played by Zetta Allas, in a wonderfully hostile manner
and those who stand accused, and the suspense rose throughout the scene
as we in the audience hoped dearly for their innocence to be proved; it
almost is before the final bit of the scene, where Rosalie is coerced
into lying, solidifying Mary's accusation. The heightened frustration of
being falsely accused was portrayed excellently by Viren Shetty, Emily
Hassan and Clara Falkowska as they, outraged, fight back against the
wonderfully dislikeable portrayal of Mary, by Cora Wilson.

In conclusion. \textit{American Classics} was a slick, well executed, well
acted anthology of great American theatre that covered a wide range of
tones and time periods to create an entertaining thought-provoking look
at many themes throughout the show. The transitions were interesting and
very well done. The acting was spot on and at times very moving or very
comedic. The accents, though occasionally imperfect, were on the whole
accurate. It was altogether a great performance, and your correspondent
sincerely enjoyed all three visits.

\end{multicols}

	\chapter{Music}
	
	
		
		On the next few pages, an excerpt from Jonathan Watt's Sonata in D is published. A full copy will be available at a later date. Details will be published in the Editor's Note, or in the Music section.
		
		
	
	\includepdf[pages={1-6}]{sonata.pdf}
	
	\includepdf[pages={2}]{cover.pdf}
	
\end{document}
